\chapter{Introduzione}
\label{chap:introduction}

Introduzione di quello di cui si parlerà nella tesi.

Contestualizzare il problema: sistemi embedded distribuiti che devono coordinarsi
in modo autonomo e resiliente, senza un orchestratore centrale, ed in un contesto
resource-constrained, e Rust come linguaggio memory-safe ed ad alte prestazioni
per la programmazione embedded.

Introduzione a YAAIR ed al concetto di Network, e Zenoh come soluzione per la
comunicazione.

\section{Contributo della tesi}
\label{sec:contribution}

Presentare in modo sintetico il lavoro svolto.

\section{Struttura della tesi}
\label{sec:thesis-structure}

Descrivere brevemente il contenuto di ciascun capitolo, una riga per
capitolo. Es.: "Il Capitolo~\ref{chap:background} introduce i concetti
di background necessari alla comprensione del contributo..."
